\documentclass{article}
\usepackage{amsmath}
\usepackage{amssymb}
\usepackage{geometry}
\usepackage{booktabs}

\geometry{a4paper, margin=1in}

\title{Assignment Solution: Electricity Consumption Prediction}
\author{}
\date{}

\begin{document}

\maketitle

\section*{Problem Statement}
We are given a dataset of 5 shops with Operating Hours ($x$) and Consumption ($y$).
The model is linear: $\hat{y} = w_0 + w_1 x$.
The loss function is Squared Error: $J(w_0, w_1) = \frac{1}{2}\sum_{i=1}^{5}(\hat{y}_i - y_i)^2$.

Dataset:
\begin{center}
\begin{tabular}{c c c}
\toprule
Shop & $x$ & $y$ \\
\midrule
1 & 1 & 3 \\
2 & 2 & 5 \\
3 & 4 & 9 \\
4 & 6 & 13 \\
5 & 8 & 17 \\
\bottomrule
\end{tabular}
\end{center}

\section*{Part (a): Loss Function \& Convexity}

\subsection*{1. Formulate the Loss Function}
Substituting the model $\hat{y}_i = w_0 + w_1 x_i$ into the loss function:
$$ J(w_0, w_1) = \frac{1}{2} \sum_{i=1}^{5} (w_0 + w_1 x_i - y_i)^2 $$

Substituting the dataset values $(x_i, y_i)$:
$$
J(w_0, w_1) = \frac{1}{2} \left[
(w_0 + 1w_1 - 3)^2 +
(w_0 + 2w_1 - 5)^2 +
(w_0 + 4w_1 - 9)^2 +
(w_0 + 6w_1 - 13)^2 +
(w_0 + 8w_1 - 17)^2
\right]
$$

\subsection*{2. Compute the Hessian Matrix}
To find the Hessian, we first determine the first-order partial derivatives (Gradient) using the Chain Rule.
Recall $J(w_0, w_1) = \frac{1}{2}\sum_{i=1}^{N}(\hat{y}_i - y_i)^2$ where $\hat{y}_i = w_0 + w_1 x_i$.

\textbf{Step 1: First Derivatives (Gradient)}
$$ \frac{\partial J}{\partial w_0} = \sum_{i=1}^{N} (\hat{y}_i - y_i) \frac{\partial \hat{y}_i}{\partial w_0} = \sum_{i=1}^{N} (w_0 + w_1 x_i - y_i) \cdot (1) $$
$$ \frac{\partial J}{\partial w_1} = \sum_{i=1}^{N} (\hat{y}_i - y_i) \frac{\partial \hat{y}_i}{\partial w_1} = \sum_{i=1}^{N} (w_0 + w_1 x_i - y_i) \cdot (x_i) $$

\textbf{Step 2: Second Derivatives (Hessian Elements)}
Differentiating the gradient components again:
1. With respect to $w_0$:
$$ \frac{\partial^2 J}{\partial w_0^2} = \frac{\partial}{\partial w_0} \sum_{i=1}^{N} (w_0 + w_1 x_i - y_i) = \sum_{i=1}^{N} 1 = N $$
Substituting $N=5$:
$$ \frac{\partial^2 J}{\partial w_0^2} = 5 $$

2. With respect to $w_0$ and $w_1$ (Mixed partial derivative):
$$ \frac{\partial^2 J}{\partial w_0 \partial w_1} = \frac{\partial}{\partial w_1} \sum_{i=1}^{N} (w_0 + w_1 x_i - y_i) = \sum_{i=1}^{N} x_i $$
Summing the $x$ values:
$$ \sum x_i = 1 + 2 + 4 + 6 + 8 = 21 $$
$$ \frac{\partial^2 J}{\partial w_0 \partial w_1} = 21 $$
(Note: By symmetry of continuous functions, $\frac{\partial^2 J}{\partial w_1 \partial w_0} = 21$).

3. With respect to $w_1$:
$$ \frac{\partial^2 J}{\partial w_1^2} = \frac{\partial}{\partial w_1} \sum_{i=1}^{N} (w_0 + w_1 x_i - y_i) x_i = \sum_{i=1}^{N} x_i \cdot x_i = \sum_{i=1}^{N} x_i^2 $$
Summing the $x^2$ values:
$$ \sum x_i^2 = 1^2 + 2^2 + 4^2 + 6^2 + 8^2 = 1 + 4 + 16 + 36 + 64 = 121 $$
$$ \frac{\partial^2 J}{\partial w_1^2} = 121 $$

Thus, the Hessian matrix $H$ is:
$$
H = \begin{bmatrix}
\frac{\partial^2 J}{\partial w_0^2} & \frac{\partial^2 J}{\partial w_0 \partial w_1} \\
\frac{\partial^2 J}{\partial w_1 \partial w_0} & \frac{\partial^2 J}{\partial w_1^2}
\end{bmatrix} = \begin{bmatrix}
5 & 21 \\
21 & 121
\end{bmatrix}
$$

\subsection*{3. Prove Convexity}
Mathematically, a continuously differentiable function of several variables is strictly convex if its Hessian matrix is \textbf{Positive Definite} for all points in its domain.

For a symmetric $2 \times 2$ matrix $H = \begin{bmatrix} A & B \\ B & C \end{bmatrix}$, the conditions for Positive Definiteness are:
\begin{enumerate}
    \item $A > 0$
    \item $\det(H) = AC - B^2 > 0$
\end{enumerate}
Alternatively, the matrix is Positive Definite if all its eigenvalues are strictly positive.

\textbf{Method 1: Sylvester's Criterion (Determinant check)}
\begin{itemize}
    \item First leading principal minor: $H_{1,1} = 5 > 0$. (Satisfied)
    \item Second leading principal minor (Determinant):
    $$ \det(H) = (5)(121) - (21)(21) = 605 - 441 = 164 $$
    Since $164 > 0$, the determinant is positive.
\end{itemize}
Since all leading principal minors are positive, $H$ is Positive Definite.

\textbf{Method 2: Eigenvalue Analysis}
We find the eigenvalues $\lambda$ by solving the characteristic equation $\det(H - \lambda I) = 0$:
$$ \det \begin{bmatrix} 5-\lambda & 21 \\ 21 & 121-\lambda \end{bmatrix} = 0 $$
$$ (5-\lambda)(121-\lambda) - 21^2 = 0 $$
Expanding the terms:
$$ (605 - 5\lambda - 121\lambda + \lambda^2) - 441 = 0 $$
$$ \lambda^2 - 126\lambda + 164 = 0 $$
Using the quadratic formula $\lambda = \frac{-b \pm \sqrt{b^2 - 4ac}}{2a}$:
$$ \lambda = \frac{126 \pm \sqrt{(-126)^2 - 4(1)(164)}}{2} $$
$$ \lambda = \frac{126 \pm \sqrt{15876 - 656}}{2} = \frac{126 \pm \sqrt{15220}}{2} $$
$$ \sqrt{15220} \approx 123.37 $$
The two eigenvalues are:
$$ \lambda_1 \approx \frac{126 + 123.37}{2} = 124.685 $$
$$ \lambda_2 \approx \frac{126 - 123.37}{2} = 1.315 $$

Since both eigenvalues $\lambda_1 > 0$ and $\lambda_2 > 0$, the Hessian matrix is Positive Definite.

\textbf{Conclusion:} Because the Hessian matrix $H$ is constant and positive definite everywhere, the Loss Function $J(w_0, w_1)$ is a \textbf{strictly convex function}.

\section*{Part (b): Gradient Calculation}
The gradient vector is given by:
$$ \nabla J = \begin{bmatrix} \frac{\partial J}{\partial w_0} \\ \frac{\partial J}{\partial w_1} \end{bmatrix} = \begin{bmatrix} \sum_{i=1}^5 (\hat{y}_i - y_i) \\ \sum_{i=1}^5 (\hat{y}_i - y_i)x_i \end{bmatrix} $$

Let $e_i = \hat{y}_i - y_i = (w_0 + w_1 x_i - y_i)$.
$$ \nabla J = \begin{bmatrix} \sum e_i \\ \sum e_i x_i \end{bmatrix} $$

\section*{Part (c): Gradient Descent Iterations}
Initial parameters: $w^{(0)} = [2, 2]^T$.
Initial residuals with $w^{(0)}$:
\begin{itemize}
    \item $i=1: \hat{y}=2+2(1)=4, y=3 \implies e=1$
    \item $i=2: \hat{y}=2+2(2)=6, y=5 \implies e=1$
    \item $i=3: \hat{y}=2+2(4)=10, y=9 \implies e=1$
    \item $i=4: \hat{y}=2+2(6)=14, y=13 \implies e=1$
    \item $i=5: \hat{y}=2+2(8)=18, y=17 \implies e=1$
\end{itemize}
Initial Gradient at $w^{(0)}$:
$$ \frac{\partial J}{\partial w_0} = \sum 1 = 5 $$
$$ \frac{\partial J}{\partial w_1} = \sum 1 \cdot x_i = 21 $$
$$ \nabla J(w^{(0)}) = \begin{bmatrix} 5 \\ 21 \end{bmatrix} $$

\subsection*{Scenario 1: Constant Learning Rate $\eta = 0.05$}

\textbf{Iteration 1:}
$$ w^{(1)} = w^{(0)} - \eta \nabla J(w^{(0)}) $$
$$ w^{(1)} = \begin{bmatrix} 2 \\ 2 \end{bmatrix} - 0.05 \begin{bmatrix} 5 \\ 21 \end{bmatrix} = \begin{bmatrix} 2 - 0.25 \\ 2 - 1.05 \end{bmatrix} = \begin{bmatrix} 1.75 \\ 0.95 \end{bmatrix} $$

\textbf{Iteration 2:}
Calculate gradient at $w^{(1)} = [1.75, 0.95]^T$.
Errors $e_i = 1.75 + 0.95 x_i - y_i$:
\begin{itemize}
    \item $x=1: 1.75 + 0.95 - 3 = -0.3$
    \item $x=2: 1.75 + 1.90 - 5 = -1.35$
    \item $x=4: 1.75 + 3.80 - 9 = -3.45$
    \item $x=6: 1.75 + 5.70 - 13 = -5.55$
    \item $x=8: 1.75 + 7.60 - 17 = -7.65$
\end{itemize}
Gradient components:
$$ \frac{\partial J}{\partial w_0} = \sum e_i = -0.3 - 1.35 - 3.45 - 5.55 - 7.65 = -18.3 $$
$$ \frac{\partial J}{\partial w_1} = \sum e_i x_i = -0.3(1) - 1.35(2) - 3.45(4) - 5.55(6) - 7.65(8) $$
$$ = -0.3 - 2.7 - 13.8 - 33.3 - 61.2 = -111.3 $$
$$ \nabla J(w^{(1)}) = \begin{bmatrix} -18.3 \\ -111.3 \end{bmatrix} $$

Update:
$$ w^{(2)} = w^{(1)} - \eta \nabla J(w^{(1)}) $$
$$ w^{(2)} = \begin{bmatrix} 1.75 \\ 0.95 \end{bmatrix} - 0.05 \begin{bmatrix} -18.3 \\ -111.3 \end{bmatrix} = \begin{bmatrix} 1.75 + 0.915 \\ 0.95 + 5.565 \end{bmatrix} = \begin{bmatrix} 2.665 \\ 6.515 \end{bmatrix} $$

\subsection*{Scenario 2: Decaying Learning Rate $\eta_t = 0.1 e^{-0.4t}$}

\textbf{Iteration 1 ($t=1$):}
$$ \eta_1 = 0.1 e^{-0.4(1)} \approx 0.1(0.67032) = 0.06703 $$
Using $\nabla J(w^{(0)}) = [5, 21]^T$:
$$ w^{(1)} = \begin{bmatrix} 2 \\ 2 \end{bmatrix} - 0.06703 \begin{bmatrix} 5 \\ 21 \end{bmatrix} = \begin{bmatrix} 2 - 0.33515 \\ 2 - 1.40763 \end{bmatrix} = \begin{bmatrix} 1.66485 \\ 0.59237 \end{bmatrix} $$

\textbf{Iteration 2 ($t=2$):}
$$ \eta_2 = 0.1 e^{-0.4(2)} = 0.1 e^{-0.8} \approx 0.1(0.44933) = 0.04493 $$
Calculate gradient at $w^{(1)} = [1.66485, 0.59237]^T$.
(Using computed values from Python script for precision):
$$ \nabla J(w^{(1)}) \approx \begin{bmatrix} -26.24 \\ -156.37 \end{bmatrix} $$
Update:
$$ w^{(2)} = \begin{bmatrix} 1.66485 \\ 0.59237 \end{bmatrix} - 0.04493 \begin{bmatrix} -26.24 \\ -156.37 \end{bmatrix} $$
$$ w^{(2)} \approx \begin{bmatrix} 1.66485 + 1.17889 \\ 0.59237 + 7.02596 \end{bmatrix} = \begin{bmatrix} 2.8437 \\ 7.6183 \end{bmatrix} $$

\section*{Part (d): Stability Analysis}
The results show that both scenarios lead to divergence (increasing loss) within the first two iterations, but Scenario 2 is \textbf{less stable} initially.

\begin{itemize}
    \item \textbf{Scenario 1 (Constant $\eta=0.05$):} The loss increased from $2.5$ to $51.6$ to $1398.1$.
    \item \textbf{Scenario 2 (Decaying):} The loss increased from $2.5$ to $101.3$ to $2135.8$.
\end{itemize}

\textbf{Comment:} Although decaying learning rates generally improve stability in the long run, in this specific case, the initial learning rate $\eta_1 \approx 0.067$ was higher than the constant rate $\eta=0.05$. Since $\eta=0.05$ was already too large (causing overshoot and divergence), the higher initial rate in Scenario 2 caused an even larger overshoot. The decay in the second step ($\eta_2 \approx 0.045$) was not sufficient to recover from the bad position reached after the first step. Therefore, the decaying schedule performed worse in these first two iterations.

\section*{Part (e): Line Search Methods}
We search along the direction $d = -\nabla J(w^{(0)}) = [-5, -21]^T$.
We minimize $\phi(\eta) = J(w^{(0)} + \eta d)$.
Initial Interval: $[0, 1]$.

\subsection*{1. Binary Search}
We compare $\phi(m)$ and $\phi(m + \epsilon)$ where $\epsilon = 10^{-3}$.

\textbf{Iteration 1:}
Current Interval: $[0, 1]$. Midpoint $m = 0.5$.
Evaluate:
$$ \phi(0.5) \approx 7006.5 $$
$$ \phi(0.501) \approx 7035.0 $$
Since $\phi(m) < \phi(m+\epsilon)$, the function is increasing at $m=0.5$. The minimum must be to the left (smaller $\eta$).
\textbf{Reduced Interval:} $[0, 0.5]$.

\textbf{Iteration 2:}
Current Interval: $[0, 0.5]$. Midpoint $m = 0.25$.
Evaluate:
$$ \phi(0.25) \approx 1695.3 $$
$$ \phi(0.251) \approx 1709.3 $$
Since $\phi(m) < \phi(m+\epsilon)$, the function is increasing at $m=0.25$. The minimum is to the left.
\textbf{Reduced Interval:} $[0, 0.25]$.

\subsection*{2. Golden-Section Search}
Initial Interval $[a, b] = [0, 1]$.
Interior points: $M_1 = 1/4 = 0.25$, $M_2 = 3/4 = 0.75$.

\textbf{Iteration 1:}
Evaluate:
$$ \phi(M_1) = \phi(0.25) \approx 1695.25 $$
$$ \phi(M_2) = \phi(0.75) \approx 15936.25 $$
Since $\phi(M_1) < \phi(M_2)$, the minimum lies in $[a, M_2]$.
\textbf{Reduced Interval:} $[0, 0.75]$.

\textbf{Iteration 2:}
New Interval $[a, b] = [0, 0.75]$. Width $= 0.75$.
Valid choice for new $M_1, M_2$ (using $1/4$ and $3/4$ of new interval):
$$ M_1 = 0 + 0.25(0.75) = 0.1875 $$
$$ M_2 = 0 + 0.75(0.75) = 0.5625 $$
Evaluate (for completeness, though not explicitly asked to reduce again):
$$ \phi(0.1875) \approx 932.8 $$
$$ \phi(0.5625) \approx 8899.7 $$
Since $\phi(M_1) < \phi(M_2)$, the new interval would be $[0, 0.5625]$.

\end{document}
